\chapter{Pythonic 慣用寫法與風格}
\label{ch:pythonic}

你大概讀過那種句子：每個字都認得，文法也挑不出毛病，但讀起來就是怪，一看就知道是從別的
語言硬翻過來的。反過來也一樣——聽一個學中文沒多久的朋友講話，每個字都用對了，組起來就是
不像我們平常講話的樣子。他沒有講錯，只是帶著口音。

程式碼也有口音。從第 \ref{ch:variables} 章到第 \ref{ch:recursion} 章，語法該教的都教完
了：你會寫變數、條件、迴圈、函式、檔案、例外、comprehension、generator、遞迴。同一件事，
用 Python 寫得出來的方法常常不只一種，而且它們\textbf{全都能動}——這一章要談的不是哪一種寫法
對、哪一種錯，是\textbf{讀的人要不要多花力氣}。

上一章的結尾，還有第 \ref{ch:iterators-generators} 章的結尾，都留了同一句話：「第
\ref{ch:pythonic} 章會把『寫得像 Python』這件事整個收起來。」就是現在。

先把一件事講在前面，免得整章讀起來像在推銷：這一章的最後一節會反過來談「什麼時候不要太
Pythonic」。腔調本身不是目的，把話講清楚才是。

這一章分三段推進：

\begin{itemize}
\item 第一段是低風險的日常慣例：程式碼該長什麼樣子、條件怎麼寫成一句話、值怎麼搬運。
這三節的共同論證只有一條——程式碼被讀的次數，遠比被寫的次數多
（第 \ref{sec:pep8}～\ref{sec:value-idioms} 節）。
\item 第二段是這一章存在的理由。名字牌與 EAFP 這兩個從第 \ref{ch:variables} 章跟第
\ref{ch:files-exceptions} 章一路跟著你走到這裡的想法，會在這裡走完一圈：從「解釋為什麼
會發生怪事的圖像」，變成「你打字的當下會用到的判準」。中間補上 \texttt{with} 這塊一直
缺著的拼圖（第 \ref{sec:is-vs-equal}～\ref{sec:eafp-closure} 節）。
\item 第三段是煞車。學會一個寫法之後，還要學會什麼時候不用它
（第 \ref{sec:not-too-pythonic} 節）。
\end{itemize}

\section{大家講好的樣子：PEP 8 與命名慣例}
\label{sec:pep8}

先看一段程式。它能跑，答案也對。

\begin{example}{python}{(tidy-todo-before.py) 能動，但看得很累}
maxitems=5
l=["Milk","Eggs","Toast"]
def addItem(item,items):
    if len(items)<maxitems:
        items.append(item)
    return items
print(addItem("Apple",l))
\end{example}

\iosample{}{['Milk', 'Eggs', 'Toast', 'Apple'] \\}

沒有一行寫錯。但你讀它的時候，眼睛要多做好幾件事：\texttt{maxitems=5} 這一串黏在一起，
你得自己在心裡把等號兩邊切開；\texttt{l} 是什麼？是 list 嗎？還是 1？\texttt{addItem}
跟 \texttt{maxitems} 一個用大寫駝峰、一個全小寫，看不出規律；\texttt{def} 前面沒有空行，
整段擠成一坨。

同樣的程式，換一種排法：

\begin{example}{python}{(tidy-todo-after.py) 一樣的程式，照慣例排過}
MAX_ITEMS = 5
shopping_list = ["Milk", "Eggs", "Toast"]


def add_item(item, items):
    if len(items) < MAX_ITEMS:
        items.append(item)
    return items


print(add_item("Apple", shopping_list))
\end{example}

\iosample{}{['Milk', 'Eggs', 'Toast', 'Apple'] \\}

輸出一模一樣。Python 完全不在乎你選哪一種——但下一個讀它的人在乎，而下一個讀它的人，
第一個是三天後的你自己。第 \ref{ch:variables} 章第 \ref{sec:naming-practices} 節就提醒過
一次：幾天後回頭看，你會完全忘記某個短名字到底是什麼。這一節要把那句話擴大到整份檔案的
外觀上。

\subsection*{PEP 8 是什麼}

Python 社群對「程式碼該長什麼樣子」有一份講好的文件，叫做 \textbf{PEP 8}。
\index{PEP 8}
PEP 是 Python Enhancement Proposal 的縮寫，直譯是「Python 改進提案」——任何人想對 Python
提一件事，都可以寫成一份 PEP 給整個社群看，通過的就變成大家的共識。PEP 8 是最老的幾份
之一，它談的不是語言功能，而是排版與命名。

有一件事要立刻講清楚，因為這一章最後一節整節都建立在它上面：

\keypoint{\textbf{PEP 8 不是語法規則。} Python 不會因為你不照它做就報錯，也不會跑得比較
慢。它是一群人為了讓彼此的程式碼看起來像同一個世界的東西，自願講好的約定。約定跟規則
不一樣的地方在於：規則由編譯器執行，約定只能靠人守。}

\subsection*{排版：四條真的有差的}

PEP 8 很長，但對現在的你真正有差的只有四條，而且每一條都有理由，不是「規定就是這樣」。

\begin{itemize}
\item \textbf{縮排用 4 個空白，而且不要跟 tab 混用。} 4 個空白是視覺上剛好看得出層次、
又還留得下空間給多層縮排的寬度。至於 tab 跟空白混用會發生什麼事，第 \ref{ch:introduction}
章第 \ref{sec:editor-pitfalls} 節示範過那個錯誤訊息了——那次的教訓在這裡變成一條慣例。
\item \textbf{一行不要太長。} PEP 8 給的數字是 79 個字元。這個數字有歷史原因，但今天更
實際的理由是：真正在工作的人常常同時開好幾個檔案並排在同一個螢幕上，一行太長，並排就
放不下了。
\item \textbf{用空行做視覺分組，但不要過量。} 兩個函式之間空兩行、函式裡面的邏輯段落之間
空一行，讀的人一眼就知道哪幾行是一組。全部黏在一起看不出結構，每兩行就空一行也一樣看不
出結構。
\item \textbf{運算子前後各留一個空白。} \texttt{x = 5} 讀起來比 \texttt{x=5} 舒服；
式子一長，差別更明顯：\texttt{len(items) < MAX\_ITEMS} 一眼看得出是在比大小，
擠成 \texttt{len(items)<MAX\_ITEMS} 就得自己找斷點在哪。理由跟中文標點一樣，
空白是在幫讀的人斷句。
\end{itemize}

\subsection*{命名：四條}

\begin{itemize}
\item \textbf{變數與函式用 \texttt{snake\_case}}：全部小寫，單字之間用底線連接。
\texttt{shopping\_list}、\texttt{add\_item} 都是這種寫法。第 \ref{ch:functions} 章教
函式命名時已經提過這個詞，這裡只是把它一起放進慣例清單。
\item \textbf{常數用全大寫加底線}，例如 \texttt{MAX\_ITEMS}、\texttt{CONFIG\_PATH}。
\index{全大寫命名慣例}
這裡要特別說清楚一件事：\textbf{Python 沒有「常數」這種東西。} 你寫成
\texttt{MAX\_ITEMS = 5}，之後照樣可以寫 \texttt{MAX\_ITEMS = 99}，Python 不會攔你。
全大寫只是一個講好的暗號，意思是「這個值是設定，不打算在程式跑的過程中被改動」。它防
不了任何事，它只是在跟讀的人講話。
\item \textbf{class 用大寫開頭的駝峰式}，例如 \texttt{ShoppingCart}。class 是第
\ref{ch:objects} 章的主題，這裡先知道有這條慣例就好。
\item \textbf{布林值用 \texttt{is\_}／\texttt{has\_} 開頭}，例如
\texttt{is\_passing}、\texttt{has\_ticket}。這一條第 \ref{ch:variables} 章第
\ref{sec:naming-practices} 節已經完整講過，不重講。
\end{itemize}

\stylepoint{這些命名慣例本身沒有一條是神聖的——真正有價值的是\textbf{整份專案裡保持
一致}。一份檔案裡一半 \texttt{snake\_case}、一半駝峰式，比全部用駝峰式還難讀，因為讀的人
會一直懷疑這個差別是不是有意義的。}

\subsection*{有工具會替你檢查這件事}

風格慣例不需要用背的。有人寫了程式來幫你抓這類問題：\texttt{flake8} 會把不合慣例的地方
一條一條列給你看，\texttt{black} 更直接，它會直接把你的檔案照慣例重排一遍。這種工具的來源
跟第 \ref{ch:data-tools} 章裝過的那些套件一樣，用 \texttt{pip} 安裝。這裡點名它們，只是
要你知道「不必用背的」，實際怎麼裝、怎麼設定，等你真的需要的時候再查就好。

\subsection*{但是}

給完規則之後，馬上要補一句「但是」——這件事很重要，重要到 PEP 8 自己就有一節在講。那一節
的標題叫「A Foolish Consistency Is the Hobgoblin of Little Minds」（愚蠢的一致性是小心眼
的妖怪），裡面寫著：

\begin{quote}
與風格指南保持一致很重要。但最重要的是：知道什麼時候該不一致——有時候風格指南就是不適用。
有疑慮的時候，用你自己最好的判斷。
\end{quote}

也就是說，這份講好程式碼該長什麼樣子的文件，自己就留了一道門，讓你在有理由的時候走出去。
什麼樣的理由算數，第 \ref{sec:not-too-pythonic} 節會整節談。

最後補一句你現在用不到、但很快會用到的事：一群人一起改同一個專案的時候，會有一道叫
code review 的關卡——你的程式碼要先被別人讀過才進得去。風格一致在那裡不是美學問題，是
直接變成別人讀你程式碼要花多少力氣。那套流程長什麼樣子，第 \ref{ch:industry} 章會講。

\begin{drills}[label=drills:pep8]
\item 把下面這段程式照這一節的慣例改寫（命名與排版都要改）：
\begin{verbatim}
TotalItems=0
def CountThem(TheList):
    global TotalItems
    for x in TheList:
        TotalItems=TotalItems+1
    return TotalItems
\end{verbatim}
\item 你寫了 \texttt{MAX\_RETRY = 3}，後來在程式的另一個地方寫了
\texttt{MAX\_RETRY = 10}。Python 會阻止你嗎？如果不會，那全大寫這個慣例到底在幫誰？
\item PEP 8 說一行不要超過 79 個字元，但你的團隊決定改用 99 個字元當上限，而且整個專案
都照這個數字。這樣算不算「錯」？
\end{drills}

\section{把條件寫成一句話：\texttt{not}、\texttt{in}、連續比較}
\label{sec:condition-idioms}

上一節談的是程式碼的外觀，這一節往裡面走一層：同一個判斷，怎麼寫會讓下一個人一眼看懂。
這一節有三個小技巧，其中兩個你手上已經有材料了，只有最後一個是全新的語法。

\subsection*{\texttt{if not items:}}

第 \ref{ch:conditionals} 章第 \ref{sec:truthy-falsy} 節教過 Python 的一個脾氣：空的東西
在 \texttt{if} 裡會被當成 \texttt{False}。那時候教的是\textbf{正面}的形式——
\texttt{if name:} 就是「name 有內容的話」。第 \ref{ch:files-exceptions} 章第
\ref{sec:eafp-recap} 節也用容器示範過一次。

現在把否定的形式補上，順便第一次把完整的清單集中列出來。下面這些值放進 \texttt{if} 裡，
一律被當成 \texttt{False}：

\begin{center}
\begin{tabular}{@{}ll@{}}
\toprule
值 & 是什麼 \\
\midrule
\texttt{0}、\texttt{0.0} & 數字的零 \\
\texttt{""} & 空字串 \\
\texttt{[]} & 空 list \\
\texttt{()} & 空 tuple \\
\texttt{\{\}} & 空 dict \\
\texttt{set()} & 空 set \\
\texttt{None} & 「沒有值」這個值 \\
\bottomrule
\end{tabular}
\end{center}

除了這些，其他一律當成 \texttt{True}。有了這張表，「這個清單是不是空的」就有三種寫法，
而且三種都能動：

\begin{example}{python}{(empty-check-three-ways.py) 三種問法，同一個答案}
items = []

if len(items) == 0:
    print("A: nothing to do")

if items == []:
    print("B: nothing to do")

if not items:
    print("C: nothing to do")
\end{example}

\iosample{}{A: nothing to do \\ B: nothing to do \\ C: nothing to do \\}

三行都印出來了，所以問題不是對錯，是哪一種讀起來比較直接。

A 版繞了一圈：先算長度、再拿長度跟 0 比。你要問的其實不是「長度是幾」，是「有沒有東西」。
B 版更繞，它造了一個空 list 出來，只為了拿去跟另一個 list 比。而且 B 版換一種容器就得
重寫——換成字串要改成 \texttt{== ""}，換成 dict 要改成 \texttt{== \{\}}。

C 版直接問「這裡面沒東西吧？」，而且它對上面那張表裡的每一種值都通用：空 list、空字串、
空 dict、空 set，寫法完全一樣。

\stylepoint{要問一個容器是不是空的，寫 \texttt{if not items:}，不要寫
\texttt{if len(items) == 0:} 或 \texttt{if items == []:}。三種都能動，但第一種不用繞道
算長度、也不用造一個空容器出來比對，而且換成任何一種容器都不必改寫。正面的形式同理：
\texttt{if items:} 就是「裡面有東西的話」。}

\subsection*{一個誠實的例外}

\texttt{if not items:} 有一個它問不出來的問題：它分不出「空的」跟「根本沒有」。

\begin{example}{python}{(empty-or-missing.py) 兩種完全不同的情況，同一個答案}
def report(items):
    if not items:
        print("nothing to do")
    else:
        print(len(items), "items")


report([])
report(None)
\end{example}

\iosample{}{nothing to do \\ nothing to do \\}

兩次都印一樣的話。但這兩件事在現實裡差很多：\texttt{[]} 是「我確實去查了，結果是一份空
清單」，\texttt{None} 是「根本沒有這份清單」——可能是檔案沒讀到，可能是函式忘了回傳東西。
一個是「今天沒有待辦事項」，一個是「待辦清單不見了」。

要分辨這兩件事，得換一個問法，而那個問法就是下一節整節在講的東西。這裡先記住這個缺口。

\subsection*{\texttt{in} 取代一長串 \texttt{or}}

第二個技巧更小。你想問「今天是不是週末」：

\begin{example}{python}{(weekend-check.py) 兩種問法}
day = "Sun"

if day == "Sat" or day == "Sun":
    print("Weekend (or version)")

if day in ("Sat", "Sun"):
    print("Weekend (in version)")
\end{example}

\iosample{}{Weekend (or version) \\ Weekend (in version) \\}

上面那一行，\texttt{day} 這個名字出現了兩次，等號出現了兩次；選項如果有五個，就會出現
五次。下面那一行，\texttt{day} 只出現一次，而且「這幾個選項」被放在一起、看起來就像一份
清單。

\texttt{in} 你早就會了——第 \ref{ch:loops} 章第 \ref{sec:break-search} 節說過「\texttt{in}
其實就是幫你把搜尋迴圈包好了」。這裡只是把同一個工具用在另一個地方。順帶一提，如果選項
很多、而且程式要反覆檢查很多次，第 \ref{ch:dict-set} 章第 \ref{sec:set-membership} 節
講過改用 set 會快很多，寫法一樣是 \texttt{in}。

\subsection*{連續比較}

這是這一章唯一一個全新的語法。你想確認一個分數在 0 到 100 之間：

\begin{example}{python}{(score-in-range.py) 夾在中間的寫法}
score = 87

if score >= 0 and score <= 100:
    print("valid (and version)")

if 0 <= score <= 100:
    print("valid (chained version)")
\end{example}

\iosample{}{valid (and version) \\ valid (chained version) \\}

Python 允許把比較運算子串起來寫，這叫\textbf{連續比較}。
\index{連續比較}
\texttt{0 <= score <= 100} 的意思就是「score 大於等於 0，而且 score 小於等於 100」，
跟上面那一行完全等價——差別是它讀起來就是數學課上寫的那個樣子，而且 \texttt{score} 這個
名字只出現一次。名字只出現一次是有實際好處的：如果你要把 \texttt{score} 改成
\texttt{final\_score}，\texttt{and} 版本有兩個地方要改，改漏一個就是一個安靜的錯。

這個寫法有一條界線要守：\textbf{只有在「同一個東西被夾在中間」的時候才用它}。

\begin{code}{python}{連續比較該用與不該用的形狀}
0 <= score <= 100          # 好：score 被夾在 0 跟 100 中間
lower <= value < upper     # 好：同一個 value 被夾住

a < b == c                 # 不要：三個不同的東西串在一起，讀的人得停下來想
\end{code}

最後一行是合法的 Python，但沒有人這樣寫。連續比較的價值來自它剛好對應到「夾在中間」這個
大家都懂的形狀，一旦串的是三個不相干的比較，那個形狀就沒了，讀的人反而要多想一次。

這一節的三個技巧看起來各做各的，其實它們在做同一件事：\textbf{把一句判斷，還原成一句
人話}。「這裡沒東西」「今天是週末之一」「分數在 0 到 100 之間」——這三句話你講給人聽的
時候本來就是這樣講的，程式碼只是照著寫下來。第 \ref{ch:conditionals} 章第
\ref{sec:readable-conditions} 節那條線（把複雜的條件拆成有名字的布林變數）講的也是同一
件事，只是切入的角度不同。

\begin{drills}[label=drills:condition-idioms]
\item 把下面這段改寫成用 \texttt{in}：
\begin{verbatim}
if ext == ".jpg" or ext == ".png" or ext == ".gif":
    print("image")
\end{verbatim}
\item 把下面這段改寫成連續比較：
\begin{verbatim}
if temperature >= 18 and temperature <= 26:
    print("comfortable")
\end{verbatim}
\item 下面四行分別會印出 \texttt{True} 還是 \texttt{False}？
\begin{verbatim}
items = []
print(not items)
print(items is None)
items = None
print(not items)
print(items is None)
\end{verbatim}
\item 什麼情況下 \texttt{if not items:} 不夠用？舉一個具體的例子。
\end{drills}

\section{值的搬運：unpacking 與 f-string}
\label{sec:value-idioms}

這一節不教任何新語法。一行都不教。

第 \ref{ch:tuples} 章第 \ref{sec:tuple-unpacking} 節已經把 unpacking 完整教完了——基本
寫法、\texttt{*rest}、用底線忽略不要的值、\texttt{a, b = b, a} 交換兩個變數，連風格要點
都給過了。第 \ref{ch:strings} 章三節也把 f-string 教完了。你手上什麼都有。

那這一節在做什麼？書上教過\textbf{怎麼寫}，沒教過\textbf{什麼時候該想起它}。這是兩件不同
的事：學會一個工具跟看到情況會反射性拿出那個工具，中間差了一段路。這一節就是那段路——
一張「看到左邊，想到右邊」的對照表。

\begin{center}
\begin{tabular}{@{}p{5.4cm}p{5.4cm}p{3.8cm}@{}}
\toprule
看到自己寫成這樣 & 該想到 & 教在哪 \\
\midrule
\texttt{temp = a} 之後 \texttt{a = b}、\texttt{b = temp} &
\texttt{a, b = b, a} &
第 \ref{ch:tuples} 章第 \ref{sec:tuple-unpacking} 節 \\
\addlinespace
\texttt{for i in range(len(items)):} 再用 \texttt{items[i]} &
\texttt{for index, item in enumerate(items):} &
第 \ref{ch:loops} 章第 \ref{sec:enumerate} 節 \\
\addlinespace
\texttt{first = row[0]} 加上 \texttt{rest = row[1:]} &
\texttt{first, *rest = row} &
第 \ref{ch:tuples} 章 \\
\addlinespace
迴圈裡拿了 key，再回頭用 \texttt{d[key]} 取值 &
\texttt{for key, value in d.items():} &
第 \ref{ch:tuples} 章 \\
\addlinespace
用不到的那個值，也認真替它取了名字 &
用底線 \texttt{\_} 接住它 &
第 \ref{ch:tuples} 章 \\
\addlinespace
\texttt{"Hello, " + name + "! You are " + str(age)} &
f-string &
第 \ref{ch:strings} 章第 \ref{sec:f-string-basics} 節 \\
\bottomrule
\end{tabular}
\end{center}

左邊那一欄，每一條都能動，而且每一條都是從別的地方帶來的口音——尤其是
\texttt{for i in range(len(items))}。很多語言的迴圈就是這樣寫的：先要一個編號，再拿編號
去取東西。Python 也允許你這樣寫，但它是全書最典型的一種腔調，因為你真正想要的是「東西
本身」，編號只是取到東西的手段。

把左右兩邊並排看一次：

\begin{code}{python}{帶著口音的版本}
temp = a
a = b
b = temp

for i in range(len(items)):
    print(i, items[i])

first = row[0]
rest = row[1:]

for key in scores:
    print(key, scores[key])
\end{code}

\begin{code}{python}{同樣的四件事，本地話}
a, b = b, a

for index, item in enumerate(items):
    print(index, item)

first, *rest = row

for key, value in scores.items():
    print(key, value)
\end{code}

實際跑一次那個迴圈的例子，確認兩種寫法真的一模一樣：

\begin{example}{python}{(list-with-index.py) 同一份清單，兩種走法}
items = ["Milk", "Eggs", "Toast"]

for i in range(len(items)):
    print(i, items[i])

print("---")

for index, item in enumerate(items):
    print(index, item)
\end{example}

\iosample{}{0 Milk \\ 1 Eggs \\ 2 Toast \\ --- \\ 0 Milk \\ 1 Eggs \\ 2 Toast \\}

\subsection*{f-string 那一條，值得單獨講}

第 \ref{ch:strings} 章介紹 f-string 的時候只說了一句「變數一多就難維護」，沒有真的把
另一種寫法擺出來給你看。現在擺出來。

用 \texttt{+} 把一句話串起來，會踩到兩個實際的痛。第一個痛會直接讓程式當掉：

\begin{example}{python}{(greeting-plus.py) 用加號串一句話，忘了轉型}
name = "Alice"
age = 30

print("Hello, " + name + "! You are " + age + " years old.")
\end{example}

\begin{verbatim}
Traceback (most recent call last):
  File "greeting-plus.py", line 4, in <module>
    print("Hello, " + name + "! You are " + age + " years old.")
          ~~~~~~~~~~~~~~~~~~~~~~~~~~~~~~~~^~~~~
TypeError: can only concatenate str (not "int") to str
\end{verbatim}

\texttt{TypeError: can only concatenate str (not "int") to str}，意思是「字串只能跟字串
相加，不能跟 int 相加」。\texttt{age} 是數字，得先用 \texttt{str()} 包起來才行。這個錯不
難修，但它會反覆發生——每次你在句子裡多塞一個數字，就得記得包一次。

第二個痛不會報錯，所以更煩：標點跟空白全部要自己數。\texttt{"Hello, "} 後面那個空白、
\texttt{"! You are "} 前後那兩個空白，少一個就是輸出擠在一起，多一個就是中間多一格，而且
它們散落在一堆引號跟加號之間，很難一眼看出來。

\begin{example}{python}{(greeting-fstring.py) 兩種寫法，同一句話}
name = "Alice"
age = 30

print("Hello, " + name + "! You are " + str(age) + " years old.")
print(f"Hello, {name}! You are {age} years old.")
\end{example}

\iosample{}{Hello, Alice! You are 30 years old. \\ Hello, Alice! You are 30 years old. \\}

下面那一行裡，那句話本來的樣子是完整的——空白跟標點就長在它們該在的位置上，你不用在腦中
把五個片段拼起來，而且數字不必轉型。

\stylepoint{這張對照表裡的寫法，沒有一個是必要的——左邊那些全部能動。它們的共同好處只有
一個：\textbf{讓「你在做什麼」直接寫在程式碼的表面上}，不用讀的人自己拼。
\texttt{a, b = b, a} 表面上就寫著「交換」，\texttt{enumerate} 表面上就寫著「我要編號跟
東西」，f-string 表面上就是那句話本來的樣子。}

\begin{drills}[label=drills:value-idioms]
\item 把下面這段改寫成用 \texttt{enumerate}：
\begin{verbatim}
names = ["Alice", "Bob", "Cathy"]
for i in range(len(names)):
    print(str(i + 1) + ". " + names[i])
\end{verbatim}
\item 把下面這段改寫成 f-string：
\begin{verbatim}
city = "Taipei"
temp = 26
print("Today in " + city + " it is " + str(temp) + " degrees.")
\end{verbatim}
\item 下面這段會出什麼錯？錯誤訊息是什麼？怎麼修？
\begin{verbatim}
count = 3
print("You have " + count + " new messages.")
\end{verbatim}
\end{drills}

\section{\texttt{is} 還是 \texttt{==}：名字牌模型的收束}
\label{sec:is-vs-equal}

上一節結尾留了一個問題：要分辨「空的」跟「根本沒有」，該怎麼寫？答案很短，但要拿到它
之前，得先把一件跟了你十七章的東西處理掉。

\subsection*{一張走過的路}

名字牌這個比喻——變數不是裝東西的盒子，是貼在物件上的名字牌——在這本書裡出場過五次：

\begin{center}
\begin{tabular}{@{}p{5.2cm}p{9.2cm}@{}}
\toprule
在哪裡 & 那一次它解釋的是什麼 \\
\midrule
第 \ref{ch:variables} 章第 \ref{sec:name-tags} 節 &
為什麼 \texttt{x = 5} 之後 \texttt{x = "hi"} 完全合法 \\
\addlinespace
第 \ref{ch:lists} 章第 \ref{sec:aliasing} 節 &
為什麼改了 \texttt{b}，\texttt{a} 也跟著變了 \\
\addlinespace
第 \ref{ch:functions} 章第 \ref{sec:function-nametag}、\ref{sec:mutable-default} 節 &
為什麼函式裡的 \texttt{append} 會影響外面，以及 \texttt{cart=[]} 那個坑 \\
\addlinespace
第 \ref{ch:tuples} 章第 \ref{sec:tuple-vs-list} 節 &
為什麼「不能改」的 tuple 裡面的 list 還是改得動 \\
\addlinespace
第 \ref{ch:advanced-functions} 章第 \ref{sec:closure} 節 &
為什麼函式回傳之後，裡面那個計數器還記得上次的數字 \\
\bottomrule
\end{tabular}
\end{center}

這張表不是要你回去重讀那五節。它是要你注意一件你大概沒發現的事：\textbf{這五次，名字牌都是被
拿來解釋「一件看起來很怪的事，為什麼會發生」}。它一直是一張事後的圖，用來收拾直覺跟
事實對不上的場面。

這一節要做的，是換掉它的身分。

\subsection*{兩個問法}

\texttt{==} 跟 \texttt{is} 問的是兩個不同的問題：

\begin{itemize}
\item \texttt{a == b} 問的是：\textbf{這兩張名字牌指的物件，裝的東西一不一樣？}
\item \texttt{a is b} 問的是：\textbf{這兩張名字牌，是不是貼在同一個物件上？}
\end{itemize}

第 \ref{ch:lists} 章第 \ref{sec:aliasing} 節那個範例最後印出 True、True、True、False，
就是這兩個問題各問兩次的結果。這裡不重貼那段程式，只用一個更短的版本確認一次：

\begin{example}{python}{(compare-carts.py) 兩台內容一樣的購物車}
my_cart = ["Milk", "Eggs"]
your_cart = ["Milk", "Eggs"]

print(my_cart == your_cart)
print(my_cart is your_cart)
\end{example}

\iosample{}{True \\ False \\}

裝的東西一樣（\texttt{==} 是 True），但它們是兩個各自造出來的 list，不是同一個
（\texttt{is} 是 False）。

\subsection*{那我日常到底什麼時候該用 \texttt{is}？}

這個問題從第 \ref{ch:lists} 章就欠著了。那裡說 \texttt{is} 用在「比較少見的情境」，卻
沒告訴你少見到什麼程度。現在給答案，而且答案短得有點意外：

\keypoint{日常寫程式，你要的幾乎永遠是 \texttt{==}。\texttt{is} 只有一個常見用途：
\textbf{跟 \texttt{None} 比}。不確定的時候，用 \texttt{==}——它通常就是你要的，而且對
\texttt{None} 也管用。}

為什麼 \texttt{None} 特別？因為 \texttt{None} 全世界只有一個。這句話第 \ref{ch:functions}
章第 \ref{sec:mutable-default} 節講過了：那一節教的安全寫法 \texttt{if cart is None:}
正是這條規則最好的既有例子。整個程式裡不管哪裡出現 \texttt{None}，都是同一個物件，所以
問「是不是同一個」跟問「值一不一樣」會得到一樣的答案，而前者更直接。

\subsection*{兩個反例}

第一個反例是 \texttt{if x == None:}。它能動，答案也對，但沒有人這樣寫。既然全世界只有
一個 \texttt{None}，「是不是同一個」才是你真正想問的問題；繞去問「值一不一樣」，是拿一
把更大的工具做一件更小的事。

\stylepoint{跟 \texttt{None} 比，寫 \texttt{if x is None:}／\texttt{if x is not None:}，
不要寫 \texttt{if x == None:}／\texttt{if x != None:}。兩種都能動，但前者是所有 Python
程式碼裡的統一寫法，讀的人掃過去不用停。}

第二個反例是真的會害人的。看到上面那條「跟 \texttt{None} 比用 \texttt{is}」之後，很容易
順手把 \texttt{is} 用到別的地方去：

\begin{example}{python}{(is-zero.py) 拿 is 去比一個數字}
count = 0

if count is 0:
    print("empty")
\end{example}

\begin{verbatim}
is-zero.py:3: SyntaxWarning: "is" with 'int' literal. Did you mean "=="?
  if count is 0:
empty
\end{verbatim}

注意兩件事。第一，Python 發出的是 \texttt{SyntaxWarning}（警告），不是錯誤——
\index{SyntaxWarning}
程式沒有停下來，它照樣印出了 \texttt{empty}。第二，訊息裡 Python 直接問你「你是不是想
寫 \texttt{==}？」，它已經看出你大概搞錯了。

這一類的錯屬於第 \ref{ch:strings} 章第 \ref{sec:f-string-basics} 節說過的「安靜的錯」：
程式跑完了、有輸出、看起來一切正常，但它成立的理由不是你以為的那個。\texttt{count is 0}
這次剛好是 \texttt{True}，是因為 Python 內部對很小的整數做了一些最佳化，讓它們共用同一
個物件——換一個大一點的數字、或換一個字串，同樣的寫法就可能變成 \texttt{False}，而且一樣
不會報錯。

這就是為什麼那條規則要寫得那麼死：\textbf{跟 \texttt{None} 比用 \texttt{is}，其他一律
\texttt{==}}。不是因為別的用法一定會錯，是因為它什麼時候會錯、什麼時候不會錯，取決於你
不需要知道、也不該依賴的細節。

\subsection*{回到上一節那個問題}

現在可以回答「空的」跟「根本沒有」怎麼分了：

\begin{example}{python}{(missing-vs-empty.py) 三種情況，三種回答}
def report(items):
    if items is None:
        print("no list at all")
    elif not items:
        print("list is empty")
    else:
        print(len(items), "items")


report(None)
report([])
report(["Milk", "Eggs"])
\end{example}

\iosample{}{no list at all \\ list is empty \\ 2 items \\}

順序很重要：先問 \texttt{is None}，再問 \texttt{not items}。反過來的話，
\texttt{not items} 會把 \texttt{None} 一起吞掉，因為 \texttt{None} 本身也是 falsy。

\subsection*{名字牌，從這裡開始換一個身分}

名字牌從第 \ref{ch:variables} 章走到這裡，一直在替你解釋「為什麼會發生怪事」。每一次它
出現，都是因為某件事跟你的直覺對不上，需要一張圖來收拾。

現在它變成一件更小、也更常用的東西：\textbf{一條你打字的時候會用到的規則。}

跟 \texttt{None} 比用 \texttt{is}，其他一律 \texttt{==}。就這樣。你不需要每次打字都想起
那張圖，也不需要在心裡重演一次名字牌貼上去的過程——這條規則已經把那張圖要講的事收在裡面
了。而如果哪一天你忘了它為什麼是這樣，那張圖還在第 \ref{ch:variables} 章那裡等你回去看。

\begin{drills}[label=drills:is-vs-equal]
\item 下面這段會印出什麼？
\begin{verbatim}
a = [1, 2, 3]
b = [1, 2, 3]
c = a
print(a == b)
print(a is b)
print(a is c)
\end{verbatim}
\item 把 \texttt{if result == None:} 改成慣用寫法。如果要問的是「result 不是 None」呢？
\item \texttt{if status is "done":} 為什麼危險？它一定會出錯嗎？
\item 你要判斷兩台購物車裝的東西一不一樣（不在乎它們是不是同一台車）。該用
\texttt{==} 還是 \texttt{is}？
\end{drills}

\section{\texttt{with} 到底在解決什麼問題}
\label{sec:with-explained}

第 \ref{ch:files-exceptions} 章第 \ref{sec:reading-files} 節第一次教開檔案的時候，除了
把 \texttt{with} 的作用講成一句話（「這段縮排區塊做完之後，自動幫我把檔案關掉」）跟一個
比喻（用完自動幫你把號碼牌繳回去、置物櫃自動關好），還留了一句話：

\begin{quote}
你之後在很多情境都會看到 \texttt{with} 這個寫法，不只用在檔案上。
\end{quote}

這一節來把那句話講完。

順帶把當時講不完的另外半句也補上：那一章第 \ref{sec:finally} 節說過「\texttt{with} 背後
的機制，精神上跟 \texttt{finally} 很像」。當時只能講到「精神上」，因為 \texttt{finally}
是那一節才剛教的。現在你兩個都有了，可以把那句話從比喻換成一段真的程式碼。

\subsection*{不用 \texttt{with} 的話，要做哪些事？}

這個問題你答得出來：開檔、讀、關檔。

\begin{example}{python}{(open-close-by-hand.py) 自己開、自己關}
file_object = open("shopping_list.txt")
content = file_object.read()
file_object.close()

print(content)
\end{example}

\iosample{}{Milk \\ Eggs \\ Toast \\ Apple \\}

三行，看起來也還好。問題是這三行有三種壞掉的方式。

\textbf{第一種：忘了寫 \texttt{close()}。} 每個開著的檔案都會佔掉一點資源，而且大部分
作業系統對「同時能開幾個檔案」有上限。開一兩個沒事，一支程式在迴圈裡開了幾千個而且都沒關，
就會出事。更常見的後果是寫檔：沒有正確關閉的檔案，最後那一段內容可能根本還沒真的寫進去。

\textbf{第二種：\texttt{close()} 放太早。} 關掉之後才發現還要再讀一次，於是對著一個已經
關掉的檔案動手。這種錯的訊息等一下就會看到。

\textbf{第三種，也是最要命的：中間那一行炸了。} 如果 \texttt{read()} 那一行丟出例外，
程式會直接跳走，\texttt{close()} 那一行根本輪不到執行。前兩種是你會忘記，第三種是你就算
沒忘記也躲不掉——因為那一行寫在那裡，只是永遠沒被跑到。

\subsection*{把 \texttt{with} 拆回 \texttt{try/finally}}

第三種問題，你其實已經有工具可以解決了。第 \ref{ch:files-exceptions} 章第
\ref{sec:finally} 節教過 \texttt{finally}：不管 \texttt{try} 裡發生什麼事，
\texttt{finally} 裡的事一定會做。所以自己開檔而想寫得安全，要寫成這樣：

\begin{example}{python}{(open-with-finally.py) 不用 with，但要寫得安全}
file_object = open("shopping_list.txt")
try:
    content = file_object.read()
finally:
    file_object.close()

print(content)
\end{example}

\iosample{}{Milk \\ Eggs \\ Toast \\ Apple \\}

現在把 \texttt{with} 版本擺在旁邊：

\begin{code}{python}{同一件事的兩種寫法}
# 手寫版
file_object = open(filename)
try:
    content = file_object.read()
finally:
    file_object.close()

# with 版
with open(filename) as file_object:
    content = file_object.read()
\end{code}

\keypoint{\texttt{with} 就是上面那五行的簡寫。它替你保證的事只有一件，但那件事很重要：
\textbf{不管縮排區塊裡發生什麼——正常跑完、中途 \texttt{return}、還是丟出例外——收尾那
一步一定會執行。} 這正是第 \ref{sec:finally} 節那個 \texttt{finally} 在做的事，只是
\texttt{with} 把它縮成一行，而且連 \texttt{close()} 都不用你自己寫。}

到這裡，「\texttt{with} 是什麼」就有一個具體的答案了：它不是魔法，它是一段 \texttt{try/
finally} 的固定寫法，被包成一個關鍵字。

\subsection*{離開 \texttt{with} 之後，那個名字還在}

有一件事值得親眼看一次，因為它常被誤會：

\begin{example}{python}{(file-after-with.py) 離開之後再去碰那個檔案}
with open("shopping_list.txt") as file_object:
    content = file_object.read()

print(file_object)
print(file_object.closed)
print(file_object.read())
\end{example}

\begin{verbatim}
<_io.TextIOWrapper name='shopping_list.txt' mode='r' encoding='UTF-8'>
True
Traceback (most recent call last):
  File "file-after-with.py", line 6, in <module>
    print(file_object.read())
          ~~~~~~~~~~~~~~~~^^
ValueError: I/O operation on closed file.
\end{verbatim}

三行輸出各講一件事。第一行：離開 \texttt{with} 之後，\texttt{file\_object} 這個名字還在，
\texttt{print} 印得出東西——\texttt{with} 不像函式那樣會建立一個獨立的空間，區塊裡的名字
出來之後照樣用得到。第二行：那個檔案物件的 \texttt{closed} 屬性是 \texttt{True}，它確實
被關掉了。第三行：真的想再讀一次，就會拿到
\texttt{ValueError: I/O operation on closed file.}（對一個已經關掉的檔案做輸入輸出）。

用上一節剛收好的那條規則的原話來說：名字牌沒有被撕掉，被改變的是它指著的那個物件的狀態。

\subsection*{「不只用在檔案上」是什麼意思}

回到那張欠條。\texttt{with} 建立的是一段\textbf{有開頭、有結尾的區間}：進去的時候做一件
事，出來的時候一定做另一件事——不管你是正常走出來，還是中途被例外叫走。

開檔跟關檔只是這種「一對」裡最常見的一組。其他常見的組合，你以後遇到的時候就會認得：

\begin{itemize}
\item 連上資料庫，用完斷開連線。
\item 打開一個壓縮檔往裡面塞東西，塞完把它封好。第 \ref{ch:automating} 章做備份的時候
就會再用到一次同樣的寫法，那時候你會發現它跟開檔案長得一模一樣。
\end{itemize}

至於「怎麼自己做一個 \texttt{with} 用得上的東西」，那需要先學會 class，超出這本書的範圍。
好消息是你幾乎不需要——標準函式庫跟常見的套件早就替你做好了，你要做的是\textbf{看到
\texttt{with} 的時候知道它在保證什麼}，以及看到一段「開了就一定要關」的程式碼時，想起有
這個東西。

寄物櫃那個比喻現在可以說得更準一點：不管你是正常辦完事走出來，還是中途被叫走，那個櫃子
都會被關好，號碼牌都會被繳回去。

\begin{drills}[label=drills:with-explained]
\item 下面這段會印出什麼？（提示：不只一行）
\begin{verbatim}
with open("notes.txt") as f:
    text = f.read()

print(f.closed)
print(f.read())
\end{verbatim}
\item 把下面這段改寫成用 \texttt{with}：
\begin{verbatim}
log_file = open("log.txt", "a")
log_file.write("done\n")
log_file.close()
\end{verbatim}
\item 如果 \texttt{with} 區塊裡的程式丟出例外，檔案還會被關上嗎？為什麼？
\end{drills}

\section{先做再說：EAFP 是一整套，不是一個語法}
\label{sec:eafp-closure}

第 \ref{ch:files-exceptions} 章第 \ref{sec:lbyl-vs-eafp} 節教過一個縮寫：EAFP，
easier to ask forgiveness than permission——與其先申請許可，不如先做了再道歉。那一章還
用第 \ref{sec:eafp-recap} 節整節回頭指出，這套邏輯遠比「檔案讀寫」這一件事大。

這一節要再往上一層。往上一層具體是三件事。

\subsection*{一、它不是一個語法，是一整個家族}

第 \ref{ch:files-exceptions} 章講的是「面對可能失敗的事，Python 偏好先做」。但這個偏好
影響的不只是你怎麼寫 \texttt{try}——它長在整個 Python 的設計裡。下面這張表裡的每一項，
你都已經在用了，只是從來沒有人告訴你它們是同一家人：

\begin{center}
\begin{tabular}{@{}p{3.6cm}p{3.6cm}p{7.2cm}@{}}
\toprule
你已經在用的東西 & 教在哪 & 它是這個家族的哪一種形狀 \\
\midrule
\texttt{try} / \texttt{except} &
第 \ref{ch:files-exceptions} 章第 \ref{sec:lbyl-vs-eafp} 節 &
本體：直接做，錯了接住 \\
\addlinespace
\texttt{dict.get(key, default)} &
第 \ref{ch:dict-set} 章第 \ref{sec:dict-get} 節 &
Python 替你把「查不到就用預設值」整包做好了 \\
\addlinespace
\texttt{set.discard()} &
第 \ref{ch:dict-set} 章第 \ref{sec:set-membership} 節 &
同上：不在裡面就當沒事，不必先問 \\
\addlinespace
\texttt{if items:} / \texttt{if not items:} &
第 \ref{ch:conditionals} 章、第 \ref{sec:condition-idioms} 節 &
語言自己把「有沒有內容」換成答案，不必你先量長度 \\
\addlinespace
\texttt{with} &
第 \ref{sec:with-explained} 節 &
\texttt{try/finally} 的簡寫——負責收尾的那一半 \\
\addlinespace
讀得懂的錯誤訊息 &
第 \ref{ch:debugging} 章 &
既然失敗是預期會發生的事，訊息當然要寫得能讀 \\
\bottomrule
\end{tabular}
\end{center}

前面三格用程式碼寫出來，你會發現它們短得幾乎不像在處理失敗：

\begin{code}{python}{同一個家族的三種說法}
value = scores.get(name, 0)    # 查不到？給 0，不必先問在不在
tags.discard("draft")          # 不在裡面？當沒事，不必先問在不在
if not items:                  # 空的？直接問，不必先量長度
    print("nothing to do")
\end{code}

這張表裡最值得停一下的是 \texttt{with} 那一格。上一節剛把 \texttt{with} 拆回
\texttt{try/finally}，所以現在你看得懂它為什麼在這張表裡：處理失敗有兩半，一半是「接住
它」（\texttt{except}），一半是「不管有沒有失敗都要收尾」（\texttt{finally}），而
\texttt{with} 就是後面那一半的固定寫法。第 \ref{ch:files-exceptions} 章沒辦法把這一格
填進來，因為那時候 \texttt{with} 還只是一個「照著寫就對了」的關鍵字。

\subsection*{二、這個偏好是從哪來的}

不是因為 Python 比較潮，也不只是因為第 \ref{ch:files-exceptions} 章講過的那個理由（檢查
完到真的動手之間，狀態可能已經變了）。往上一層的說法是這樣的：

\keypoint{Python 把「失敗」當成一件\textbf{會正常發生的事}，而不是一件\textbf{該被預防的
意外}。所以它給你的是接住失敗的工具，而不是要求你事先把所有可能出錯的情況都想過一遍。
查不到、開不起來、轉不了型——這些事本來就會發生，程式的工作是處理它們，不是假裝它們不
存在。}

這句話解釋了很多你早就見過、卻沒想過為什麼的事。為什麼 \texttt{dict} 有一個
\texttt{get()}？為什麼 \texttt{set} 除了 \texttt{remove()} 還要有一個 \texttt{discard()}？
為什麼 Python 的錯誤訊息會長到用一整段話告訴你哪裡不對，還畫箭頭指出是哪個位置——如果
失敗被當成不該發生的意外，訊息只要寫「失敗」兩個字就夠了。

\subsection*{三、兩條線在這裡會合}

這本書從頭到尾只有兩個心智模型：名字牌，還有 EAFP。它們一個在第 \ref{ch:variables} 章
起頭，一個在第 \ref{ch:files-exceptions} 章起頭，一路上看起來完全沒有關係——
一個在講變數怎麼運作，一個在講錯誤怎麼處理。

把它們各自要修正的那個直覺並排看，關係就出來了：

\begin{center}
\begin{tabular}{@{}p{3.8cm}p{4.4cm}p{6.2cm}@{}}
\toprule
你原本的直覺 & Python 實際做的事 & 這個直覺假設了什麼 \\
\midrule
\texttt{b = a} 會複製一份 &
只是多貼一張名字牌 &
假設 Python 會替我多做一步（複製） \\
\addlinespace
動手之前要先檢查 &
直接做，失敗了告訴你 &
假設 Python 應該替我多做一步（防呆） \\
\bottomrule
\end{tabular}
\end{center}

兩個直覺其實是同一個直覺：\textbf{以為 Python 會在你沒開口的時候，替你多做一步。}

它不會。Python 的預設是「不多做」——不替你複製，也不替你預檢。這件事的代價是你得知道它
預設不做什麼，這也是為什麼這兩件事要花十幾章反覆提醒。好處是另一面：\textbf{沒有任何一
步是在你背後偷偷發生的。} 你想要複製，就自己寫 \texttt{.copy()}；你想要檢查，就自己寫
\texttt{if}。你寫幾步，Python 就做幾步。

這兩個模型從第 \ref{ch:variables} 章跟第 \ref{ch:files-exceptions} 章分頭走了十幾章，走
到這裡才會合。

\begin{drills}[label=drills:eafp-closure]
\item 下面四段程式，哪幾段屬於這一節那張表裡的家族？
\begin{verbatim}
# (1)
name = people.get(user_id, "unknown")

# (2)
if x % 2 == 0:
    print("even")

# (3)
with open("data.txt") as f:
    text = f.read()

# (4)
if os.path.exists(path):
    f = open(path)
\end{verbatim}
\item \texttt{with} 為什麼算 EAFP 家族的一員？用一句話說明。
\item 「這個數是不是偶數」這種判斷，該不該包成 \texttt{try/except}？為什麼？
\end{drills}

\section{什麼時候不要太 Pythonic}
\label{sec:not-too-pythonic}

從第 \ref{ch:advanced-functions} 章到現在，你連續學了五章新寫法：\texttt{lambda}、
\texttt{sorted(key=)}、decorator、comprehension、generator、遞迴，再加上這一章的六節
慣用寫法。這一節是煞車。

\subsection*{同一件事，三個版本}

給你一份成績資料（假設是從檔案讀進來的），要做三件事：篩掉不及格的、換算成等第、按分數
由高到低排序。

A 版，全部擠成一行：

\begin{example}{python}{(grades-one-liner.py) A 版：一行寫完}
rows = [
    {"name": "Alice", "score": 88},
    {"name": "Bob", "score": 54},
    {"name": "Cathy", "score": 92},
    {"name": "David", "score": 71},
]

result = [(r["name"], "A" if r["score"] >= 90 else "B" if r["score"] >= 80 else "C") for r in sorted([r for r in rows if r["score"] >= 60], key=lambda r: r["score"], reverse=True)]

print(result)
\end{example}

\iosample{}{[('Cathy', 'A'), ('Alice', 'B'), ('David', 'C')] \\}

B 版，拆成幾個有名字的步驟：

\begin{example}{python}{(grades-steps.py) B 版：一步一個名字}
rows = [
    {"name": "Alice", "score": 88},
    {"name": "Bob", "score": 54},
    {"name": "Cathy", "score": 92},
    {"name": "David", "score": 71},
]


def grade_of(score):
    if score >= 90:
        return "A"
    if score >= 80:
        return "B"
    return "C"


passed = [row for row in rows if row["score"] >= 60]
ranked = sorted(passed, key=lambda row: row["score"], reverse=True)
result = [(row["name"], grade_of(row["score"])) for row in ranked]

print(result)
\end{example}

\iosample{}{[('Cathy', 'A'), ('Alice', 'B'), ('David', 'C')] \\}

C 版，完全不用這五章教過的任何東西，全部手寫迴圈：

\begin{example}{python}{(grades-plain-loop.py) C 版：全部手寫}
rows = [
    {"name": "Alice", "score": 88},
    {"name": "Bob", "score": 54},
    {"name": "Cathy", "score": 92},
    {"name": "David", "score": 71},
]

passed = []
for row in rows:
    if row["score"] >= 60:
        passed.append((row["score"], row["name"]))

passed.sort(reverse=True)

result = []
for score, name in passed:
    if score >= 90:
        grade = "A"
    elif score >= 80:
        grade = "B"
    else:
        grade = "C"
    result.append((name, grade))

print(result)
\end{example}

\iosample{}{[('Cathy', 'A'), ('Alice', 'B'), ('David', 'C')] \\}

三個版本輸出一模一樣。哪一個比較好？

先不要急著回答。

\subsection*{三條判準，而且判準的對象是人}

\begin{itemize}
\item \textbf{讀的人是誰？} 只有你自己看、跑一次就丟的小腳本，A 版沒問題——它確實比較快
寫完。要交給別人維護，或者三個月後你自己還要回來改，那就是 B 版。這個問題沒有一般答案，
因為它問的不是程式碼，是情況。
\item \textbf{一行讀不完，就沒有優點了。} 慣用寫法的價值來自「一眼看懂」，一旦要折行、
要往右捲、要從中間往兩邊讀，那個價值就不見了，剩下的只有難度。這條就是第
\ref{ch:comprehensions} 章第 \ref{sec:comp-limits} 節那條「超過兩層或兩行就拆開」的
一般化——那裡講的是 comprehension，這裡對這五章教過的每一種工具都成立。
\item \textbf{寫的時候用光了聰明，除錯的時候就沒有聰明可以用了。} 這句話改寫自
Kernighan 的一句老話：「每個人都知道，除錯比一開始寫程式難兩倍。所以，
如果你寫的時候就用盡了自己所有的聰明，你要拿什麼去除錯它？」第 \ref{ch:debugging} 章教的每一種除錯手法——讀 traceback、插
\texttt{print}、一步一步縮小範圍——全部都有一個前提：你得看得懂那段程式碼在做什麼。
\end{itemize}

\stylepoint{要不要用一個比較精簡的寫法，問三句話就夠了：\textbf{誰會讀它？一行讀不讀得
完？出錯的時候我還追得動嗎？} 三句話問的都不是程式碼本身，是人——寫程式碼的你、讀程式碼
的人、還有半年後回來修它的那個人。}

\subsection*{回答那個問題}

B 版。

理由不是「B 比較 Pythonic」——這種說法沒有意義，因為 A 版用到的 Python 特性其實比 B 版
還多。理由是很具體的一件事：A 版把篩選、排序、換算三件事壓在同一行裡，所以當結果不對的
時候，你分不出是哪一件出的錯。你沒辦法在中間插一個 \texttt{print} 看看篩完剩下幾筆，因為
「篩完」這個中間狀態根本沒有名字。B 版有 \texttt{passed}、有 \texttt{ranked}，每一步都
停得下來、看得到。

順帶一提，A 版還有一個實務上的麻煩：它有 152 個字元，在一般的螢幕寬度下讀不完。

\subsection*{反過來的錯誤，也是錯的}

讀到這裡有一個很自然的反應：那我乾脆全部寫 C 版好了，反正也能動，而且絕對不會太花俏。

這也是錯的，而且錯得沒有比較少。有兩件事要講清楚。

\textbf{第一，讀跟寫的標準不一樣。} 別人的程式碼裡一定會出現 comprehension、
\texttt{with}、\texttt{if not items:}、\texttt{a, b = b, a}——這些東西用得非常多，多到
你不可能繞開。\textbf{看得懂它們是不能談判的}；輪到自己寫的時候，你才有選擇權。第
\ref{ch:advanced-functions} 章第 \ref{sec:map-filter} 節對 \texttt{map} 跟 \texttt{filter}
說過「看得懂就好」，那句話的一般化就是這條：能讀是底線，怎麼寫是選擇。

\textbf{第二，回到這一章開頭那個比喻。} 怕講得有腔調，就乾脆不講話——那不叫沒有口音，
那叫沒有在講話。C 版不是「比較保守的選擇」，它是把一件三句話講得完的事講成十五句。

所以這一節不是在勸你少用工具，是在勸你為每一次使用給得出理由。「因為看起來比較厲害」
不是理由，「因為這樣一眼看得出我在篩什麼」是。

\subsection*{最後一句}

這一章從頭到尾沒有教你任何新東西——嚴格說起來只有一個連續比較，而那也只是把
\texttt{and} 少寫一次。它做的只有一件事：把「這樣寫可以」變成「這樣寫，下一個人讀起來
會比較輕鬆」。

誰是下一個人，只有你知道。

\begin{drills}[label=drills:not-too-pythonic]
\item 下面這一行出現在一支你要交給同事接手的程式裡。該不該拆開？
\begin{verbatim}
out = "\n".join(f"{k}: {sum(v)/len(v):.1f}" for k, v in data.items())
\end{verbatim}
\item 同事的程式碼裡有一段你看不懂的 comprehension。你該叫他改掉嗎？
\item 你寫了一支只有自己會用、跑一次就刪掉的小腳本。需要照 PEP 8 排版嗎？
\end{drills}

下一章換一個方向。到目前為止，你的資料一直是散的：一份 list 裝分數、另一份 dict 裝名字，
要用的時候自己記得它們是一組。第 \ref{ch:objects} 章要教一種把資料跟操作它的動作綁在
一起的做法，叫做 class。

\section*{練習解答}
\addcontentsline{toc}{section}{練習解答}

\paragraph{第~\ref{sec:pep8}~節}
\begin{enumerate}
\item 一種寫法：
\begin{verbatim}
total_items = 0


def count_them(the_list):
    global total_items
    for item in the_list:
        total_items = total_items + 1
    return total_items
\end{verbatim}
改了四件事：變數與函式名改成 \texttt{snake\_case}；\texttt{x} 換成看得出意思的
\texttt{item}；等號前後補空白；\texttt{def} 前面空兩行。（\texttt{TotalItems} 不是常數，
它會被改動，所以不該用全大寫。）
\item 不會阻止你，第二次指派會直接生效。全大寫幫的是\textbf{讀的人}——它在說「這個值是
設定，不打算被改」。Python 沒有常數這種東西，這個慣例防不了任何事，它只是一個講好的暗號。
\item 不算錯。PEP 8 自己就寫著「知道什麼時候該不一致」，而且行長那條特別常被團隊改。真正
會出問題的是團隊裡一半的人用 79、一半的人用 99——不一致的成本比數字本身高得多。
\end{enumerate}

\paragraph{第~\ref{sec:condition-idioms}~節}
\begin{enumerate}
\item \texttt{if ext in (".jpg", ".png", ".gif"):}
\item \texttt{if 18 <= temperature <= 26:}
\item 依序是 \texttt{True}、\texttt{False}、\texttt{True}、\texttt{True}。前兩行：空 list
是 falsy，所以 \texttt{not items} 是 \texttt{True}，但它不是 \texttt{None}。後兩行：
\texttt{None} 也是 falsy，所以 \texttt{not items} 一樣是 \texttt{True}，而這次
\texttt{is None} 也是 \texttt{True}。關鍵是第一行跟第三行答案一樣——這正是
\texttt{if not items:} 分不出「空的」跟「根本沒有」的地方。
\item 當你需要分辨「空的」跟「根本沒有」的時候。例如一個函式回傳「搜尋結果」，
\texttt{[]} 代表「找過了，沒有符合的」，\texttt{None} 代表「根本沒搜尋成功」——這兩件事
要給使用者的訊息完全不同，用 \texttt{if not items:} 會把它們混成一種。這種時候要先問
\texttt{if items is None:}。
\end{enumerate}

\paragraph{第~\ref{sec:value-idioms}~節}
\begin{enumerate}
\item
\begin{verbatim}
names = ["Alice", "Bob", "Cathy"]
for index, name in enumerate(names, start=1):
    print(f"{index}. {name}")
\end{verbatim}
（不用 \texttt{start=1} 也可以，改寫成 \texttt{f"\{index + 1\}. \{name\}"}。順帶把
\texttt{+} 串接也一起換掉了。）
\item \texttt{print(f"Today in \{city\} it is \{temp\} degrees.")}
\item 錯誤是
\begin{verbatim}
TypeError: can only concatenate str (not "int") to str
\end{verbatim}
因為 \texttt{count} 是整數，不能直接用 \texttt{+} 接在字串後面。修法有兩種：包成
\texttt{str(count)}，或直接改寫成
\texttt{print(f"You have \{count\} new messages.")}——後者不必轉型，也少數幾個空白。
\end{enumerate}

\paragraph{第~\ref{sec:is-vs-equal}~節}
\begin{enumerate}
\item 依序是 \texttt{True}、\texttt{False}、\texttt{True}。\texttt{a} 跟 \texttt{b}
內容一樣但是兩個各自造出來的 list，所以 \texttt{==} 是 \texttt{True}、\texttt{is} 是
\texttt{False}；\texttt{c = a} 只是多貼一張名字牌，所以 \texttt{a is c} 是 \texttt{True}。
\item 改成 \texttt{if result is None:}。要問「不是 \texttt{None}」的話，
慣用寫法是 \texttt{if result is not None:}，不是
\texttt{if not result is None:}——後者雖然也能動，但讀起來要多繞一次。
\item 因為它問的是「這兩個是不是同一個物件」，而你想問的其實是「內容一不一樣」。
它不一定會出錯——如果 \texttt{status} 是程式裡直接寫死的 \texttt{"done"}，Python 可能
讓兩邊共用同一個字串物件，結果剛好是 \texttt{True}；但如果 \texttt{status} 是拼出來的
（例如從檔案讀進來、或用 \texttt{join()} 組出來的），同樣的比較就會變成 \texttt{False}，
而且不會報錯。「有時候對、有時候錯，還不告訴你」比「一定錯」更難抓。正確寫法是
\texttt{if status == "done":}。
\item \texttt{==}。你要問的是內容，不是身分。用 \texttt{is} 的話，兩台裝著一樣東西的
購物車會得到 \texttt{False}——除非它們本來就是同一台。
\end{enumerate}

\paragraph{第~\ref{sec:with-explained}~節}
\begin{enumerate}
\item 先印出 \texttt{True}（檔案已經被 \texttt{with} 關掉了），然後第二個 \texttt{print}
會丟出 \texttt{ValueError: I/O operation on closed file.}，程式停在那裡。重點是
\texttt{f} 這個名字本身還在、還用得到，被關掉的是它指著的那個檔案物件。
\item
\begin{verbatim}
with open("log.txt", "a") as log_file:
    log_file.write("done\n")
\end{verbatim}
\item 會。這正是 \texttt{with} 存在的理由。把它拆回 \texttt{try/finally} 就看得很清楚：
關檔那一步寫在 \texttt{finally} 裡，而 \texttt{finally} 不管 \texttt{try} 裡發生什麼事
都會執行。如果自己寫 \texttt{open()} 加 \texttt{close()} 而中間丟了例外，\texttt{close()}
那一行根本輪不到跑。
\end{enumerate}

\paragraph{第~\ref{sec:eafp-closure}~節}
\begin{enumerate}
\item (1)、(3) 是。(1) 是 \texttt{dict.get()}，把「查不到就用預設值」包好了；(3) 是
\texttt{with}，負責收尾的那一半。(2) 不是——它是一個完全可預期、跟外部無關的條件判斷，
本來就該用 \texttt{if}，跟「可能失敗」無關。(4) 剛好是相反的那一種：它先檢查再動手，
是典型的 LBYL，而且檢查完到真的 \texttt{open()} 之間，檔案還可能被刪掉。
\item 因為它是 \texttt{try/finally} 的簡寫，保證「不管中間成不成功，收尾一定會做」——
而「假設事情可能失敗，並且先把善後安排好」正是這個家族的共同精神。
\item 不該。這是一個完全可預期、只跟手上這個數字有關的判斷，不是「可能因為外部因素失敗」
的操作，用 \texttt{if x \% 2 == 0:} 就好。EAFP 不是「所有地方都改成 \texttt{try}」，
它針對的是檔案、網路、使用者輸入、字典查詢這一類本來就可能落空的事。
\end{enumerate}

\paragraph{第~\ref{sec:not-too-pythonic}~節}
\begin{enumerate}
\item 該拆。它同時做了三件事：走過每一組資料、算平均並格式化到小數點後一位、把結果
組成多行字串，而且一行讀不完。至少拆成兩步（先算好一份 \texttt{averages}，再組字串），
出錯的時候才有地方插 \texttt{print} 看中間結果。如果這是你自己跑一次就丟的腳本，
留著也無妨——但題目說了要交給同事接手。
\item 不該叫他改掉，先去讀懂它。這是「讀跟寫的標準不一樣」那條界線：看得懂別人的慣用寫法
是不能談判的，那是你的功課不是他的。讀懂之後如果發現它真的違反了上面那三條判準（一行讀
不完、三件事壓在一起），那時候再提，而且要講的是具體理由，不是「我看不懂」。
\item 不需要，但排一下對你自己有好處。PEP 8 自己就說了「知道什麼時候該不一致」，一支馬上
要刪掉的腳本沒有下一個讀者。不過有兩件事值得留：變數名取得看得懂，還有等號前後留空白——
因為「跑一次就刪掉」的腳本，往往不會真的被刪掉。
\end{enumerate}
